\solchapter{continuum}

\sol{continuum}{1} $\sigma = F/A = 981\un{N}/(7.9\times10^{-7}\un{m^2}) = 1.25\un{GPa}$, five times the yield strength of mild steel: it would yield and fail. Piano wire survives, stretching by $\Delta L = \sigma L/E = (1.25\times10^9)(2)/(2\times10^{11}) = 1.25\un{cm}$.

\sol{continuum}{2} Isothermal: $\sqrt{P/\rho} = \sqrt{1.013\times10^5/1.20} = 290\un{m/s}$. Adiabatic: $\sqrt{\gamma P/\rho} = 343\un{m/s}$. Helium ($\gamma = 5/3$, $m = 4\un{u}$): $\sqrt{\gamma\kB T/m}\approx1000\un{m/s}$. The vocal cords set the fundamental, which is unchanged, but the resonances of the vocal tract (the formants, which give a voice its timbre) scale with the sound speed and rise about threefold: the voice sounds higher and thinner.

\sol{continuum}{3} One violin is ten times quieter: $70 - 10 = 60\un{dB}$. Reaching $90\un{dB}$ needs $100$ times the intensity of $70\un{dB}$, so $1000$ violins.

\sol{continuum}{4} Open at both ends: $f_n = nv/2L = 172$, $343$, $515\un{Hz}$. Closed at one end: $f = (2n-1)v/4L = 86$, $257$, $429\un{Hz}$.

\sol{continuum}{5} $f = \frac{v}{2\pi}\sqrt{A/V\ell} = \frac{343}{2\pi}\sqrt{(3\times10^{-4})/(7.5\times10^{-4}\times0.08)}\approx120\un{Hz}$, about the B below the bass clef (end corrections lengthen the effective neck and lower it a little). Drinking increases the air volume $V$, so the pitch falls as $V^{-1/2}$.

\sol{continuum}{6} $E\sim U_0/a^3 = (3\times1.6\times10^{-19}\un{J})/(1.5\times10^{-10}\un{m})^3\approx1.4\times10^{11}\un{Pa} = 140\un{GPa}$. Diamond is about seven times stiffer --- within an order of magnitude, which is all the estimate claims. (Diamond's dense network of short bonds, each atom bonded four ways, adds a geometric factor.)

\sol{continuum}{7} $(E/\rho g)^{1/3} = (10^{10}/(700\times9.8))^{1/3}\approx113\un{m^{1/3}}$, so with $D = 1\un{m}$, $h_{\max}\approx0.8\times113\approx90\un{m}$. With a safety factor of four, a real tree with a $1\un{m}$ trunk stands about $25\un{m}$ tall --- about right. The tallest redwoods, about $100\un{m}$, have trunks several metres across.

\sol{continuum}{8} Weight $\times1000$; bone cross-section $\times100$; stress $\times10$. To keep the stress the same, the cross-section must grow $1000$-fold, so the bones must be $\sqrt{10}\approx3$ times thicker than mere scaling gives. The giant would have legs like an elephant's, mostly bone, and would still struggle to move.

\sol{continuum}{9} $\Delta t = d(1/v_S - 1/v_P) = d(0.119\un{s/km})$, so $d\approx67\un{km}$. One station gives only the distance: a circle on the map. Three stations pin down the epicentre; four also fix the depth and the time of the quake.

\sol{continuum}{10} The energy per thickness as a function of crack half-length is $U(a) = 4a\gamma - \pi a^2\sigma^2/E$. It increases for small $a$ and decreases for large $a$, with a maximum at $\dd U/\dd a = 0$: $a_c = 2\gamma E/\pi\sigma^2$. A crack longer than $a_c$ lowers its energy by growing, so it runs. Equivalently, the fracture stress for a crack of half-length $a$ is $\sigma_f = \sqrt{2E\gamma/\pi a}$. For glass with $a = 1\un{\mu m}$: $\sqrt{2(7\times10^{10})(1)/\pi\times10^{-6}}\approx2\times10^{8}\un{Pa} = 200\un{MPa}$, against a theoretical strength of $E/10 = 7\un{GPa}$. (Window glass, with larger flaws, breaks at about $50\un{MPa}$.)

\sol{continuum}{11} A square membrane of side $L$ has modes $f = (c/2L)\sqrt{m^2+n^2}$ with $m,n\ge1$. Those with $f<F$ are the lattice points with positive coordinates inside a circle of radius $R = 2LF/c$: about a quarter of its area, $\tfrac14\pi R^2 = \pi L^2F^2/c^2 = \pi AF^2/c^2$. Weyl showed the leading term is the same for any shape, with a correction proportional to the perimeter. So counting the frequencies below $F$ for large $F$ determines the area: a perfect ear could hear it (and the perimeter too).

\sol{continuum}{12} Measurements on orchestral timpani (Fletcher and Rossing) put the $(1,1)$, $(2,1)$, $(3,1)$, and $(4,1)$ modes at about $1:1.50:1.98:2.44$ times the $(1,1)$ frequency --- close to $2:3:4:5$, the second to fifth harmonics of a tone an octave below. Because the strongest partials line up nearly as a harmonic series, the ear fuses them into one definite pitch, heard at the frequency of the principal $(1,1)$ mode. The symmetric $(0,n)$ modes move air in and out of the kettle as a whole, so they radiate efficiently and die within a fraction of a second; and the drum is struck about a quarter of the way in from the rim, where they are excited only weakly. They contribute the initial thud, not the pitch.
